\documentclass[12pt]{article}
\usepackage[paperwidth=6in,paperheight=9in]{geometry}
\geometry{
  inner=15mm,
  outer=10mm,
  top=13mm,
  bottom=13mm,
  headheight=12pt,
  headsep=7mm,
  footskip=10mm,
  includehead=true,
  includefoot=true
}
\usepackage{fontspec}
\setmainfont{Libertinus Serif}
\usepackage{microtype}
\usepackage{setspace}
\usepackage{ragged2e}

\title{Chapter 12: The Unexpected Interlocutor}
\author{From \textit{Stolen Words}}
\date{}

\raggedbottom
\clubpenalty=150
\widowpenalty=150
\setlength{\parskip}{4pt plus 2pt minus 1pt}
\setlength{\parindent}{1em}

\begin{document}

\maketitle

\vspace{1cm}

{\centering\itshape In which a conversation occurs across digital mediums\\and the architecture of authority\\reveals itself through silence.\par}
\vspace{0.5cm}

\justifying

At three seventeen in the morning—Maya would remember the exact time because her laptop's clock glowed blue against the darkness of her apartment, and she had been staring at it for the past forty minutes—an email arrived that she had not expected to receive.

The apartment, if we must describe it (and perhaps we must, for physical spaces shape intellectual decisions in ways philosophers rarely acknowledge), occupied the second floor of a building constructed in 1924, back when landlords believed in high ceilings and landlords' children believed in subdividing their inheritance. Three rooms, technically, though one served primarily as a repository for books that had overflowed the shelves in the other two. The desk where Maya sat faced a window overlooking an alley where, at this hour, nothing moved except occasional wind-borne newspapers and one persistent cat whose nighttime habits Maya had come to know better than she knew the habits of most humans.

Around her: seventeen books open to various pages. She had developed the habit—peculiar but effective—of creating what she called "conversation circles," arranging texts so that disparate authors could speak to one another across centuries. Tonight's circle included Prabhupāda's 1972 Bhagavad-gītā (the physical copy, spine cracked, pages annotated in three colors of ink), the 1983 revision (borrowed from a temple, pristine, smelling of that particular mustiness that comes from books shelved but not read), three volumes of Sanskrit commentary (Śaṅkara, Rāmānuja, and a modern critical edition whose editor had footnoted himself into incomprehensibility), two books on translation theory (one brilliant, one tedious), and—somewhat incongruously—a volume of Jorge Luis Borges essays that had nothing to do with Vaiṣṇava theology but which Maya found herself reading between bouts of textual comparison, as if Borges's labyrinths might provide relief from the tangle she had discovered in sacred transmission.

The email's sender: Devananda Swami, whose name Maya recognized immediately. Fifty years in the tradition—a prominent ISKCON guru, one of the most influential figures in the institution A.C. Bhaktivedanta Swami Prabhupāda had founded. Author of twelve books. Thousands of initiated disciples across four continents—not institutional exaggeration but documented reality, visible in the worldwide network of students who quoted his lectures, attended his seminars, and regarded him as one of the movement's leading scholars. He had studied in Vṛndāvana (Krishna's childhood home, a major pilgrimage site), taught in Māyāpur (ISKCON's spiritual headquarters in West Bengal), established temples in three European cities and two American ones, and served for decades in senior editorial positions within the Bhaktivedanta Book Trust. His photograph on his books showed a man whose face had settled into that particular expression of serene authority that comes from decades of being right, or at least of never being contradicted.

The email itself was brief:

"Ms. Rodriguez,

I have heard of your investigation into textual differences between editions of the Bhagavad-gītā As It Is. I am attaching an audio message addressing this matter. If you wish to continue this discussion, I will consider your response.

Devananda Swami"

The attachment: a single audio file hosted on a server whose URL suggested institutional infrastructure—secure, password-protected, traceable.

The audio itself proved interesting in ways that had nothing to do with its content. Devananda Swami had recorded his response rather than writing it—a choice that Maya, who had spent considerable time studying how different mediums shape different kinds of truth-claiming, found revealing. Audio permits certain rhetorical moves that text does not: the strategic pause, the sigh of exasperation, the slight elevation of voice that suggests patience tried. It also, crucially, resists the kind of close analysis that written words invite. One cannot underline a sigh. One cannot footnote a pause.

Maya downloaded the file (8.3 megabytes, MP3 format, recorded—according to the metadata—on a device whose microphone cost more than her laptop), opened her audio editing software (not to edit, but to annotate, timestamp, create what amounted to a critical edition of spoken words), and listened.

"Good morning." His voice: measured, accented with what Maya recognized as upper-caste North Indian English, the kind that signals education at institutions where Sanskrit and philosophy were taught alongside cricket and colonial administration. "I find this topic to be one that has been discussed millions of times—" here a slight laugh, not quite derisive but not quite generous either "—and exhausted. I am very familiar with the accusation that two versions lead to different paths." Pause. Three seconds. The sound of papers shuffling. "Which is absurd. Not to mention..." another pause, shorter, "...stupid."

Maya rewound. Listened again. The progression from "absurd" to "stupid" was interesting. Escalation disguised as clarification. She made a note.

"It seems to me that the people who talk like this have their own selfish personal motives." The phrase "selfish personal motives" delivered with the precise diction of someone who has used it before, often, in contexts where it effectively ended discussion. "I know the people who have strongly protested about this. I've seen the different versions. I worked for many years in the BBT—" here emphasis, the kind that invites the listener to recognize authority "—Prabhupāda trusted me to produce his books."

Maya paused the recording. \textit{Argumentum ad verecundiam}, the old name for it. Appeal to authority. She wondered if the Swami knew he was deploying classical rhetorical strategies, or if institutional life had taught him these moves so thoroughly that he performed them unconsciously, the way one learns to swim or ride a bicycle—not through theoretical understanding but through repeated submersion in the element that requires them.

"If you can give me a practical, solid example—" the words "practical" and "solid" given extra weight "—of a change like that stupid robot says in the recording you sent me..." Maya had sent no recording. She made a note of this. The Swami was responding not to her email but to some other conversation, some other critic, the amalgamated voice of all who had questioned. She had become, already, not an individual correspondent but a representative of a category: "these critics."

"...that the two versions lead to different spiritual paths, \textit{really?} Give me a practical example of that, and if it has merit, I'll accept it."

The sentence ended with the kind of finality that does not actually invite response. It was the finality of the master permitting the student to demonstrate competence before the assembly, knowing that the demonstration will fail because the criteria for success have been defined by the master and remain, necessarily, undefined for the student.

Maya responded with three documented examples, each chosen for its clarity:

\begin{enumerate}
\item The transformation of "forgotten soul" to "forgetful soul" in Bhagavad-gītā 2.13—shifting spiritual tragedy from divine to human responsibility.

\item The systematic replacement of "The Blessed Lord said" with "The Supreme Personality of Godhead said" in twenty-two instances—reframing the divine relationship from blessing-bestower to ontological superior.

\item The alteration of "all surrender" to "them surrender" in 4.11—transforming universal reciprocation into conditional response.
\end{enumerate}

The Swami's response came as another audio file, this time defensive: "These things are so honestly childish. 'Blessed Lord' versus 'Supreme Personality of Godhead'? Different tone, yes. But not a philosophical transformation. I oppose the changes, but to exaggerate that it changes everything—please, that's not for adults."

The pattern was clear: acknowledge the difference, minimize its significance, dismiss examination as exaggeration. Maya recognized it as the standard institutional defense—not denial of facts, but compartmentalization. Yes, changes exist. No, they don't matter. Yes, they're wrong. No, examining them closely is immature.

She sought a second opinion from Dr. Rāmānuja Shastri, a former ISKCON scholar now teaching in Kerala who had earned a reputation for unflinching textual analysis. His response, characteristically dense but illuminating, cut through the Swami's rhetorical moves:

"The Swami claims 'Blessed Lord' and 'Supreme Personality of Godhead' are functionally equivalent because both refer to Krishna. This assumes denotation exhausts meaning. But words aren't merely pointers—they're consciousness-shaping tools. 'Blessed' belongs to the semantic field of gift-giving, grace freely given. 'Supreme Personality of Godhead' belongs to systematic theology—it invites analysis, not approach.

"The Swami says intimacy requires \textit{gopi}-level \textit{rasa} or it's not intimate at all. This is a false dichotomy. Intimacy exists on a spectrum. 'Blessed Lord' occupies significant middle ground between pure theological distance and maximum devotional intimacy. To change from grace-centered blessing to ontological superiority is to eliminate relational warmth, not merely reduce it.

"On 'forgotten' versus 'forgetful': the Swami argues from systematic theology—Krishna never forgets, therefore souls are forgetful. But Prabhupāda wrote from the conditioned soul's experiential perspective: lost, abandoned, crying 'God has forgotten me.' One evokes divine mercy for the lost. The other assigns fault for carelessness. Different spiritualities entirely.

"The Swami's middle position—opposing changes while denying they matter—is contradiction, not moderation. Either textual changes shape consciousness or they don't. If they don't matter, why oppose them? If they do matter, why dismiss examination as childish? This reveals institutional self-preservation: fifty years of identity built on certain assumptions. To examine them too closely would require rebuilding that identity."

Maya realized the Swami's response wasn't actually addressed to her. It was an internal debate he'd been having for years, perhaps decades.

The medieval scholars had a term for this: \textit{duplex veritas}—double truth. The capacity to hold contradictory positions simultaneously by assigning them to different domains. One could recognize something evidentially while denying it practically.

Maya understood what she had gained from the exchange. Not persuasion—she had never expected that—but confirmation. The institutional response to documented evidence would be: acknowledge but minimize, oppose but defend, recognize but compartmentalize. This wasn't unique to the Swami. This was structural, the ordinary way institutions preserve themselves against evidence that threatens foundational narratives.

If the defense was structural rather than personal, the solution couldn't be personal either. It would require documentation so thorough that compartmentalization became impossible, evidence so systematic that minimization failed.

The conversation with the Swami was over. The investigation had merely begun.

\end{document}
